\section{Discussion}
\label{sec:discussion}

\subsection{What does the V-shape tell us about LLM computation?}

The striking observation that the V-axis at layer 14 is nearly orthogonal
(cosine 0.067) to the V-axis at layer 28, while layer 27 is 95\% aligned
with layer 28, suggests a specific computational story.

Early layers (L1--L3) inherit a \emph{word-level} valence signal from the
token embeddings. Word embeddings like ``happy'' and ``sad'' are already
far apart in embedding space (this is well known from
word2vec-era work). Probing at L1 recovers this signal at AUC 0.75 on SST-2.

Middle layers (L4--L26) perform context integration. During this phase, the
word-level valence signal is mixed with syntactic, discourse, and task-relevant
features, and the \emph{linear} valence direction is destroyed. A PCA probe
at any middle layer recovers essentially nothing (AUC $\approx 0.53$).

Final layers (L27--L28) prepare the residual stream for the LM head's
next-token prediction. For causal language modeling, the LM head must
produce coherent next-word distributions, which requires the residual stream
to expose ``decodable'' semantic directions. We hypothesize that the final
layers perform a kind of \emph{re-linearization} of task-relevant features,
and that valence is one such feature that the LM head reads off as a nearly
linear direction.

This interpretation makes a testable prediction: non-causal models
(masked-language models or bidirectional models) should show a different
V-shape. We leave this for future work.

\subsection{Why does arousal fail?}

The valence/arousal asymmetry is the most striking empirical finding of this
paper. The same PCA-of-centroids method that gives $|r|=0.50$ for valence
gives $|r| \leq 0.16$ for arousal, across 5 model families and every test
we ran. Three hypotheses:

\begin{enumerate}[leftmargin=*]
\item \textbf{Arousal is distributed.} Perhaps arousal is encoded in a
non-linear or multi-dimensional manner in LLM residual streams, so a single
linear direction cannot capture it. A probing study with non-linear probes
(e.g., MLP) could test this.
\item \textbf{Arousal is embedding-only.} Perhaps arousal information sits
in the embedding layer (where ``calm'' and ``excited'' are far apart) but is
not re-encoded in later layers because causal language modeling does not
need it. An L1-only probe could test this.
\item \textbf{Arousal definition mismatch.} Perhaps the PCA second component
is simply capturing a different dimension (e.g., topic, subject, register)
that happens to be more variance-heavy than arousal after valence is removed.
\end{enumerate}

We find hypothesis 1 most likely, but further work is needed.

\subsection{When does \ours{} work and when does it not?}

Based on our results, \ours works when:
\begin{itemize}[leftmargin=*]
\item the target concept has a \emph{clean linear manifestation} in the
final layer of the LLM (valence yes, arousal no, toxicity no, formality
yes on self-LOOCV but untested on external data);
\item the downstream task rewards \emph{ranking} rather than
absolute classification accuracy (lexicon correlation, binary sentiment
with median threshold, ordinal regression);
\item inference speed matters (\speedup{} faster than direct prompting);
\item interpretability matters (a single linear direction is easier to
inspect than a prompting workflow).
\end{itemize}

It does \emph{not} work for:
\begin{itemize}[leftmargin=*]
\item raw classification accuracy (direct prompting beats us by
15\,pp on SST-2);
\item concepts that are not linearly encoded (toxicity, arousal);
\item fine-grained distinctions within a class (9-class GoEmotions nearest
centroid gives 28\%, below the majority baseline of 39\%).
\end{itemize}

\subsection{Data-Efficiency Implications}

The n-shot ablation shows that \emph{one} story per class already produces a
signal above chance (AUC 0.74 on SST-2). Fifteen stories per class match
the accuracy of a supervised classifier fit on 5{,}000 labeled samples.
This is a $\sim$30$\times$ data-efficiency gain over supervised methods.
For applications in low-resource languages, domain-specific sentiment
(legal, medical, financial), or rapid prototyping, this efficiency is the
main practical benefit of \ours.

\subsection{Theoretical Connection to Few-Shot Learning}

\ours is a degenerate case of few-shot prototype learning
\cite{snell2017prototypical} where (i)~the prototypes are LLM-generated rather
than human-labeled, (ii)~the query representation is the LLM's own final
layer, (iii)~the classifier is a single projection (not a soft-max over
distances). The key insight is that the LLM itself can serve as both the
data source and the representation source; no human labels are needed for
either role.
